\section{Zadanie 15 (Filip Manijak)}
Niech $X_1,\dots,X_n$ będzie ciągiem niezależnych zmiennych losowych
z rozkładu wykładniczego z parametrem $\lambda=1$.
Niech $\bar X$ będzie średnią arytmetyczną tego ciągu.
Jak duże musi być $n$, aby
$$
P(0.9 \le \bar X \le 1.1) \ge 0.95 \; ?
$$

\subsection*{Rozwiązanie}
Zmienne $X_i$ mają $E[X_i] = 1, Var[X_i] =1$. Wtedy $S = \frac{\sum X_i}{n}$. Skoro $n$ jest wystraczająco duże, to podstawowe CTG daje nam:
$$
P \left( a \leq \frac{S-1}{\frac{1}{\sqrt{n}}} \leq b \right) = \Phi(b)-\Phi(a)
$$
Ponieważ średnia ma $E[S] = 1, Var[S] = 1/n$.
Przekształcając:
$$
P \left( a \leq \frac{S-1}{\frac{1}{\sqrt{n}}} \leq b \right) = P \left( \frac{a}{\sqrt{n}} +1 \leq S \leq \frac{b}{\sqrt{n}}+1 \right)
$$
Czyli $a \geq -0.1 \sqrt{n}$ oraz $b \leq 0.1 \sqrt{n}$. Minimalne $n$ otrzymamy oczywiście dla najmniejszego $a$ i największego $b$, czyli $ - a = b = 0.1 \sqrt{n} $. Czyli chcemy znaleźć minimalne $n$ że:
$$
\Phi(b) - \Phi(a) = 2*\Phi(b)-1 = 2 * \Phi(0.1 \sqrt{n}) -1 \geq 0.95
$$
Niestety nie da się tego ładnie policzyć, więc sprawdzamy w tablicach. Wychodzi nam, że $\Phi(x) \geq 0.975$ dla $x \geq 1.959964$. Czyli 
$$
0.1 \sqrt{n} \geq 1.959964 \implies n \geq 385
$$

\subsection*{Odpowiedź}
\[
\boxed{\text{385}}
\]
